\documentclass[12pt,a4paper]{article}
\usepackage{natbib}
\usepackage[british]{babel}
\usepackage{parskip}
\usepackage{amsmath}
\usepackage{graphicx}
\usepackage{hyperref}
\usepackage{url}
%\usepackage{appendix}

%Any global settings go here
\graphicspath{{Figures/}}

\bibpunct[:]{(}{)}{;}{a}{,}{,}

\renewcommand{\baselinestretch}{1.5}

%some shortcut commands
\newcommand{\bi}{\begin{itemize}}
\newcommand{\ei}{\end{itemize}}
\newcommand{\be}{\begin{enumerate}}
\newcommand{\ee}{\end{enumerate}}

\begin{document}

\begin{titlepage}

\title{Investigating the effectiveness of diversification strategies based on alternative risk measures\\
\line(1,0){250}\\
\large Proposal}

\date{\today}

\author{Richard Montgomery (MNTRIC006)\\
Tayla Radmore (RDMTAY001)\\
Thomas K{\"o}nigkr{\"a}mer (KNGTHO005)}

\maketitle

\begin{figure}

\begin{center}
\includegraphics[width=3cm]{UCTlogo.jpg}\\[1cm]
\end{center}

\end{figure}

\begin{center}
\line(1,0){150}\\
{\tt mntric006@myuct.ac.za}\\
{\tt rdmtay001@myuct.ac.za}\\
{\tt thomasekng@gmail.com}
\end{center}

\thispagestyle{empty}

\end{titlepage}

\pagenumbering{roman}

\tableofcontents

\newpage

\pagenumbering{arabic}

\section{Introduction}
\label{sec:Intro}

Asset allocation is an important aspect of creating a portfolio of assets. This process is commonly referred to as the diversification strategy. The majority of literature on the matter, however, focuses specifically on diversification strategies which use volatility (i.e. variance) as the risk measure. This is illustrated in the works done by \cite{chow2011survey}, and \cite{CHEN20111179}. This practice has persevered mainly due to the computational ease of using variance as a risk measure. However, in recent times the use of variance as a risk measure has come under considerable scrutiny. In his paper \textit{Risk is not Variance: An Expository Note}, \cite{cox2008risk} mathematically proves that mean-variance decision making violates the principles of a rational decision maker. This has led to a need for research into the use of alternative risk measures in diversification strategies. With this project, we intend to assess the efficacy of diversification strategies based on select alternative risk measures.

\section{Research Question}
\label{sec:ResQues}

\subsection{Primary Research Question}
\label{subsec:PriRQ}

\textit{Are diversification strategies based on alternative risk measures more effective than those based on volitility in the South African context?}

To answer this research question, we will have to answer the secondary research questions outlined below.

\subsection{Secondary Research Questions}
\label{subsec:SecRQ}

\bi
\item What constitutes a ``good'' alternative risk measure?
\item Which alternative risk measures are best suited to our needs?
\item What diversification strategy should be used with each of the chosen risk measures?
\item What is the theoretical nature of each of these diversified portfolios under their respective alternative risk measures?
\item How do these diversified portfolios compare to modern volatility diversification methods when back-tested on historical South African share data?
\ei
 
\section{Literature Review}
\label{sec:LitRev}

\subsection{Core Literature}
\label{subsec:CorLit}

In answering the questions we have set ourselves, we must gather information on diversification methods that can be used on various different risk measures. Dowd and Blake (2006), introduced us to a number of alternative risk measures that we could consider – whilst also critiquing them and offering their own recommendations. Bruder et al (2012) is our core literature with regards to the diversification strategies which we will be making use of.

\cite{dowd2006after} focus on the broad class of quantile-based risk measures (QBRM), including primarily Value-at-Risk (VaR), coherent risk measures, spectral risk measures, and distortion risk measures.

VaR is simply defined as the maximum loss of a portfolio for a given confidence interval, and is widely used in industry \citep{consiglirisk}.  Despite this, \cite{dowd2006after} believe that VaR is an imperfect measure. In addition to this, they justify replacing it with an alternative risk measure, because doing so would require very little extra work and would result in a more sound measure of risk. Their main criticisms of VaR is that it gives no information beyond the worst case scenario that it describes. \cite{dowd2006after} go further in saying that this lack of information in the tail loss region reflects an unrealistic risk-seeking behaviour in the investor. \cite{ACERBI20021505} further supports the view of \cite{dowd2006after}, by stating that VaR does not meet the axioms of coherence. These risk measure axioms were postulated by \cite{artzner1999coherent} in order to ensure that risk is more effectively managed. We will investigate these axioms more closely in the final literature review.

\cite{dowd2006after} do concede that VaR displays some positive aspects. This paragraph will expand on these points they made. VaR can be used to compare portfolios that are not restricted to a certain type of asset. Furthermore, it takes all risk factors that affect the portfolio into account, meaning that no simplifying assumptions are needed and that the factors do not have to be added incrementally to the risk measure to avoid any complications. Finally, \cite{dowd2006after} identify the probability linked units of measure of VaR as intuitive to understood.

\cite{dowd2006after} suggest an alternative to VaR in the form of coherent risk measures, of which Expected Shortfall (ES) is a special case \citep{ACERBI20021505}. ES is defined as the expected loss beyond the maximum loss defined by VaR \citep{consiglirisk}. \cite{dowd2006after} describe it as an easily generated risk measure. They do, however, draw attention to the fact it suggests the investor is risk-neutral over the lower tail region it describes. This could prove to be problematic, since investors are generally attributed with risk-averse appetites for risk.

\cite{dowd2006after} further discuss spectral risk measures, which combine the properties of coherence risk measures with risk-aversion theory. They highlight a number of issues with risk aversion theory, which presently we shall not delve much deeper into. However, it is worthwhile to keep in mind that risk-aversion theory is not a perfect or universally accepted idea. Spectral risk measures would require us to choose a risk-aversion function for the investor. This is a subjective decision that is difficult to make and would prove challenging for the research we are doing.

Distortion measures are the last risk measures examined by \cite{dowd2006after}. They focus especially on a generalisation of the Wang Transform. They believe that it is a particularly useful measure due to its ability to recover the Capital Asset Pricing Model (CAPM) as well as the Black-Scholes model. They further describe it as a measure superior to both VaR and ES, since it takes the lower tail region of losses into account. Further investigation into this risk measure should be considered.

In our hunt for the best risk measures, it is worth mentioning the ``best practices'' risk measures, briefly discussed by \cite{dowd2006after}. They described the work done by \cite{dhaene2003economic}, which argues that there is no universally ``best'' risk measure, since one needs to take the circumstances of the particular scenario into consideration.

Risk measures only comprise a portion of our research. The difficult task will be finding a suitable diversification strategy to use given the alternative risk measures we will be using. \cite{bruder2012managing} describe risk budgeting (RB) in their paper, which they describe as an ``alternative indexing strategy''. The most important part of their paper, is the single sentence in which they state that this diversification strategy can be generalised so that it can be used for a number of different risk measures, provided the risk measure is both convex, and satisfies the Euler decomposition. \cite{bruder2012managing} did not explore this train of thought any further, but rather performed back tests comparing the returns of indices which make use of different diversification techniques. The methods and results described may prove to be useful when we begin our own back testing. 


\subsection{Further Reading}
\label{subsec:FurRead}

In addition to the works discussed, further ideas should be explored. Firstly, this includes delving into the aforementioned risk measures more deeply. Perhaps the use of other risk measures should be considered as well. Secondly, diversification methods to be used will have to be investigated more comprehensively. This, again, will either entail expanding on our understanding of RB diversification methods, or either discovering other viable methods. 

\cite{dowd2006after} mentioned a number of other risk measures. Thus, they have provided us with a list of other possible risk measures which may be worthwhile to study more deeply:

\bi
\item Convex risk measures
\item Dynamic risk measures
\item Comonotonicity approaches
\item Markov bounds approaches
\item ``Best practices'' risk measures
\ei
 
``Best practice'' risk measures (which were also mentioned in section \ref{subsec:CorLit} above) may be particularly interesting to look into during the literature review of this paper. These risk measures focus on the idea that the context of the problem may be the most important factor to consider \citep{dowd2006after}. This may, however, be beyond the scope our research question.

A particular paper that is worth drawing attention to is  \textit{Conditional Value-at-Risk, spectral risk measures and (non-)diversification in portfolio selection problems - A comparison with mean–variance analysis}, by Mario Brandtner (2013). This paper is useful in that it seems to have a similar approach to what we have in mind.

Further useful readings which we may wish to consult, in addition to those already cited, include: 
\bi
\item  \textit{The Properties of Equally Weighted Risk Contributions Portfolios}, by Maillard S., Roncalli T. and Teiletche J. 
\item \textit{Toward Maximum Diversification}, by Choueifaty Y. and Coignard Y.
\item \textit{Risk Budgeting: A New Approach to Investing}, by Rahl L.
\item \textit{Active Portfolio Management: A Quantitative Approach for Producing Superior Returns and Controlling Risk}, by Grinold R. and Kahn R.
\item \textit{Risk and Asset Allocation}, by Meucci A.
\item \textit{Challenges and standards in integrating surveys of structural variation.}, by Scherer S.W. et al.
\item \textit{The Markowitz Optimization Enigma: Is ‘Optimized’ Optimal?}, by Michaud R.O.
\item \textit{Coherent Measures of Risk}, by Artzner P. et al.
\item \textit{Coherent Representations of Subjective Risk Aversion}, by Acerbi C. 
\item  \textit{On Law Invariant Coherent Risk Measures}, by Kusuoka S.
\item  \textit{Optimal Rules for Ordering Uncertain Prospects}, by Bawa V.S.
\item \textit{Mean-Risk Analysis With Risk Associated With Below-Target Returns}, by Fishburn P.C.
\item  \textit{A Class of Distortion Operators for Pricing Financial and Insurance Risks}, by Wang S.S.
\item  \textit{Assessing VaR Accuracy}, by Dowd K.
\item  \textit{Measuring Market Risk}, by Dowd K. 
\item \textit{Bayesian Value at Risk: From Linear to Non-Linear Portfolios}, by Siu T.K. et al.
\ei

\section{Data}
\label{sec:Data}

Fortunately the necessary data are not difficult to obtain. We will need historical equities' price data over the last twenty years, say. These will be South African data, due to the focus of our paper. Our supervisor, Rowan Douglas, will help/guide us to acquire these data either from industry directly, or otherwise using the resources available to us in the library, such as the Bloomberg terminal. 

\section{Methodology}
\label{sec:Method}

First and foremost, we must devise diversification strategies using various alternative risk measures. We intend to focus on either three or four risk measures. Despite its flaws (discussed in section \ref{subsec:CorLit} above) VaR has been identified as an important risk measure to consider, due to its prominent use in industry today \citep{consiglirisk}. Since ES uses similar workings to VaR, it may be easy to include this risk measure in our research as well. Additionally, we would like to look more deeply into at least one other risk measure which seems unusual, yet useful and efficient, to us.

Given our current level of research, the RB diversification method seems to be the most feasible diversification method at our disposal. Of course, this means that we will have to ensure that the risk measures we use are compatible with this diversification technique (i.e. that they are convex and satisfy the Euler decomposition \citep{bruder2012managing}). Further research and readings will help in the specifics of using this approach, and may also bring other alternative diversification techniques to light.

We plan to perform a back test on the collected past data to assess the efficacy of the diversification strategies based on the alternative risk measures we shall be investigating. This will be compared against a benchmark: namely the diversification strategy based on volatility as its risk measure. If any similar investigations are found in the course of our research, further comparisons can be made.  The programming language used will be R.

\section{Paper Structure}
\label{sec:struc}

\be
\item Introduction
\item Background \& Theory
	\be
	\item Alternative Risk Measures
	\item Diversification Methods
	\item Back Testing Theory 
	\ee
\item Methodology \& Data
\item Results
\item Discussion \& Conclusion
\item References
\item Appendices 
\ee
 
\section{Division of Work}
\label{sec:dow} 

This is a group project. Collaborative work allows for opportunities of synergy following from a greater scope of knowledge and skills. It does, however, come with its own challenges. Richard Montgomery will be taking charge of the theoretical aspect of the project, such as the literature review. Tayla Radmore will be heading the practical, coding aspect of the project. Thomas K{\"o}nigkr{\"a}mer will act as an intermediary for these two aspects of the project: assisting in research, writing or coding, editing writing and workings, and ensuring cohesion in the project.

\section{Time-line}
\label{sec:time}

\bi
\item \textbf{Monday, 30 April 2018:} Any individual work for the literature review is to be completed by this date. Up until the next deadline, we will bring our individual parts together into a cohesive whole.
\item \textbf{Friday, 4 May 2018:} Soft deadline for literature review - to be sent to our supervisor. Over the weekend, each group member will individually inspect the draft.
\item \textbf{Monday, 7 May 2018:} Group editing of draft after individual editing over the weekend. If necessary, a meeting will be set up with our supervisor as soon as possible.
\item \textbf{Friday, 11 May 2018:} Submit literature review
\item \textbf{Pre-June Exams:} Utilise this time to:
	\bi
	\item Collect the necessary stocks' data
	\item Start work on the diversification methods to be used
	\item Possibly perform a back test, using the historical data, for the diversification strategy using volatility as its benchmark. This result will be used as a benchmark against which the other strategies will be compared. 
	\ei
\item \textbf{June/July Vac:} Utilise this time to:
	\bi
	\item Perform the back tests for the alternative strategies
	\item Assess how these strategies performed in comparison to the MPT strategy.
	\ei
\item \textbf{Monday, 23 July 2018:} Any individual work for the draft is to be completed by this date. Up until the next deadline, we will bring our individual parts together into a cohesive whole.
\item \textbf{Friday, 27 July 2018:} Soft deadline for draft paper - to be sent to our supervisor. Over the weekend, each group member will individually inspect the draft.
\item \textbf{Monday, 30 July 2018:} Group editing of draft after individual editing over the weekend. If necessary, a meeting will be set up with our supervisor as soon as possible.
\item \textbf{Friday, 10 August 2018:} Submit draft paper 
\item \textbf{Friday, 7 September 2018:} Feedback for submitted draft. The next month will be used to improve our project, which will hopefully only require polishing.
\item \textbf{Friday, 19 October 2018:} Submit final paper 
\ei

\newpage
\bibliographystyle{natbib}
\bibliography{Bibliography}
\label{bib:bibliography}
 
\end{document}
